
\documentclass[11pt]{article}

\usepackage[margin=1in]{geometry}
\usepackage{amsmath}
\usepackage{amssymb}
\usepackage{booktabs}
\usepackage{enumitem}
\usepackage{listings}
\usepackage{xcolor}
\usepackage{hyperref}

\lstset{
    basicstyle=\ttfamily\small,
    breaklines=true,
    frame=single,
    columns=fullflexible,
    showstringspaces=false
}

\title{Photon-Particle Lighting Implementation}
\author{}
\date{}

\begin{document}

\maketitle

\section{Purpose}

The current simulation represents light using photon particles emitted from a
source. Each photon travels at a fixed speed, may pass through other photons,
and interacts only with simulation boundaries.

Rendering the photons directly does not produce satisfactory illumination
because the photon distribution is too sparse. Individual photons appear as
isolated points or streaks rather than producing a continuous illuminated
surface.

The lighting system should therefore separate the following two functions:

\begin{enumerate}
    \item Photon particles transport light through the simulation.
    \item Boundary surfaces accumulate and display light received from photon
    impacts.
\end{enumerate}

Photon visibility and boundary illumination must be treated as separate
rendering calculations.

\section{Required Lighting Model}

The recommended system is based on boundary light deposition.

The required sequence is:

\begin{enumerate}
    \item Emit photon particles from the light source.
    \item Move each photon at its prescribed fixed speed.
    \item Detect photon collisions with boundaries.
    \item Determine the boundary position and boundary normal at the collision.
    \item Calculate the amount of light deposited by the photon.
    \item Spread the deposited light over a local boundary region.
    \item Accumulate the light contribution in a boundary-light buffer.
    \item Apply temporal persistence or averaging to reduce noise.
    \item Render the boundary using the accumulated illumination.
    \item Optionally render a subset of photon particles as visual tracers.
\end{enumerate}

The photon particles are therefore transport samples. They do not need to be
dense enough to visually fill the illuminated region.

\section{Separation of Photon Visibility and Surface Illumination}

The existing photon brightness calculation is based on how closely the photon
direction points toward the eye. This is appropriate only when deciding how
visible a photon tracer should be.

It is not the correct calculation for determining how much light the photon
deposits onto a boundary.

\subsection{Photon Tracer Visibility}

For rendering an optional visible photon tracer, define

\begin{equation}
    \mathbf{v}_{e}
    =
    \frac{\mathbf{x}_{e}-\mathbf{x}_{p}}
    {\left\|\mathbf{x}_{e}-\mathbf{x}_{p}\right\|},
\end{equation}

where

\begin{itemize}
    \item $\mathbf{x}_{e}$ is the eye or camera position,
    \item $\mathbf{x}_{p}$ is the photon position,
    \item $\mathbf{d}_{p}$ is the normalized photon travel direction.
\end{itemize}

The photon tracer brightness may be calculated as

\begin{equation}
    I_{\mathrm{tracer}}
    =
    I_{\min}
    +
    \left(1-I_{\min}\right)
    \left[
        \max\left(
            \mathbf{d}_{p}\cdot\mathbf{v}_{e},
            0
        \right)
    \right]^{q},
\end{equation}

where $q$ controls the angular concentration of the visible tracer.

This value affects only the appearance of the photon particle itself.

\subsection{Boundary Illumination}

When the photon strikes a boundary, the deposited illumination should depend on
the incidence angle between the incoming photon direction and the outward
boundary normal.

Let

\begin{itemize}
    \item $\mathbf{d}_{p}$ be the normalized photon direction of travel,
    \item $\mathbf{n}_{b}$ be the outward boundary normal,
    \item $E_{p}$ be the photon energy or intensity weight.
\end{itemize}

The incidence factor is

\begin{equation}
    C_{\mathrm{incidence}}
    =
    \max\left(
        -\mathbf{d}_{p}\cdot\mathbf{n}_{b},
        0
    \right).
\end{equation}

The deposited light is then

\begin{equation}
    E_{\mathrm{deposit}}
    =
    E_{p}C_{\mathrm{incidence}}.
\end{equation}

This calculation is independent of the camera position.

\section{Required Data Structures}

\subsection{Photon Structure}

Each photon requires at least the following state:

\begin{lstlisting}[language=C++]
struct PhotonParticle
{
    glm::vec4 position;      // xyz = world position
    glm::vec4 direction;     // xyz = normalized travel direction
    glm::vec4 colorEnergy;   // rgb = photon color, w = photon energy
    uint32_t active;         // 1 = active, 0 = inactive
    uint32_t emitterIndex;
    uint32_t generation;
    uint32_t padding;
};
\end{lstlisting}

The direction should remain normalized. The photon velocity can be reconstructed
using

\begin{equation}
    \mathbf{v}_{p}
    =
    c_{p}\mathbf{d}_{p},
\end{equation}

where $c_{p}$ is the fixed simulation photon speed.

If every photon uses the same speed, the speed does not need to be stored in
each photon structure.

\subsection{Boundary-Light Structure}

Each renderable boundary element requires a stored illumination value.

A boundary element may be one of the following:

\begin{itemize}
    \item a boundary particle,
    \item a boundary vertex,
    \item a surface cell,
    \item a triangle,
    \item an analytic boundary sample.
\end{itemize}

A general boundary-light structure is:

\begin{lstlisting}[language=C++]
struct BoundaryLight
{
    glm::vec4 accumulatedLight; // rgb = accumulated illumination
    glm::vec4 filteredLight;    // rgb = temporally filtered illumination
};
\end{lstlisting}

If atomic floating-point vector operations are not available, the channels may
be stored separately:

\begin{lstlisting}[language=C++]
struct BoundaryLightAtomic
{
    uint32_t red;
    uint32_t green;
    uint32_t blue;
    uint32_t hitCount;
};
\end{lstlisting}

The integer values may represent scaled fixed-point illumination.

For example,

\begin{equation}
    I_{\mathrm{integer}}
    =
    \operatorname{round}
    \left(
        S_{\mathrm{light}}I_{\mathrm{float}}
    \right),
\end{equation}

where $S_{\mathrm{light}}$ is a fixed scaling factor.

\subsection{Photon-Collision Result}

The photon-boundary collision stage should produce sufficient information for
the deposition stage.

\begin{lstlisting}[language=C++]
struct PhotonBoundaryHit
{
    uint32_t photonIndex;
    uint32_t boundaryIndex;

    glm::vec4 hitPosition;     // xyz = world-space collision position
    glm::vec4 boundaryNormal;  // xyz = outward normalized normal

    float incidence;
    float depositedEnergy;
    uint32_t valid;
    uint32_t padding;
};
\end{lstlisting}

This structure may be stored explicitly, or the deposition may be performed
immediately inside the collision shader.

Immediate deposition avoids the need for an intermediate hit buffer, but an
explicit hit buffer may be useful during debugging.

\section{GPU Buffers}

The implementation should provide the following GPU buffers.

\begin{enumerate}
    \item Photon particle buffer.
    \item Boundary geometry or boundary-particle buffer.
    \item Current-frame boundary-light deposition buffer.
    \item Temporally filtered boundary-light buffer.
    \item Optional photon-boundary hit buffer.
    \item Optional debug counter buffer.
\end{enumerate}

A possible Vulkan buffer organization is:

\begin{lstlisting}[language=C++]
VkBuffer photonBuffer;
VkBuffer boundaryBuffer;

VkBuffer boundaryLightCurrentBuffer;
VkBuffer boundaryLightFilteredBuffer;

VkBuffer photonHitBuffer;       // optional
VkBuffer photonDebugBuffer;     // optional
\end{lstlisting}

The current light buffer is cleared before each deposition pass. The filtered
light buffer persists across frames.

\section{Photon Motion}

Photon motion should remain independent of lighting accumulation.

For each active photon,

\begin{equation}
    \mathbf{x}_{p}^{n+1}
    =
    \mathbf{x}_{p}^{n}
    +
    c_{p}\mathbf{d}_{p}\Delta t.
\end{equation}

Photons may pass through other photons. Therefore, no photon-photon collision
test is required.

Only boundary intersections must be detected.

Depending on the current simulation design, boundary collision detection may
use:

\begin{itemize}
    \item boundary particles,
    \item analytic boundary equations,
    \item grid occupancy,
    \item triangle intersections,
    \item signed-distance functions.
\end{itemize}

The lighting implementation should use the same collision result already used
to reflect, absorb, or terminate the photon.

\section{Boundary Interaction}

At a boundary collision, the implementation must determine the desired photon
behavior.

The supported behaviors may include:

\begin{enumerate}
    \item absorption,
    \item reflection,
    \item partial absorption and reflection,
    \item photon termination,
    \item photon rebirth at the emitter.
\end{enumerate}

A boundary material may contain:

\begin{lstlisting}[language=C++]
struct BoundaryOpticalMaterial
{
    glm::vec4 baseColor;   // rgb = surface color
    float absorption;     // fraction absorbed
    float reflection;     // fraction reflected
    float roughness;
    float padding;
};
\end{lstlisting}

The deposited energy may be calculated as

\begin{equation}
    E_{\mathrm{deposit}}
    =
    E_{p}
    C_{\mathrm{incidence}}
    A_{b},
\end{equation}

where $A_{b}$ is the material absorption factor.

For a fully absorbing boundary,

\begin{equation}
    A_{b}=1.
\end{equation}

For a partially reflective boundary,

\begin{equation}
    0 < A_{b} < 1.
\end{equation}

The remaining photon energy is

\begin{equation}
    E_{p}^{\mathrm{new}}
    =
    E_{p}
    \left(
        1-A_{b}
    \right).
\end{equation}

\section{Spatial Light Deposition}

Depositing the entire photon contribution into one boundary sample will produce
visible speckling. Each photon hit should therefore distribute its energy over
a neighborhood of boundary samples.

Let $\mathbf{x}_{h}$ be the photon hit position and $\mathbf{x}_{j}$ be the
position of neighboring boundary sample $j$.

Define

\begin{equation}
    r_{j}
    =
    \left\|
        \mathbf{x}_{j}
        -
        \mathbf{x}_{h}
    \right\|.
\end{equation}

Only samples within an influence radius $R_{L}$ receive light.

A compact quadratic kernel is recommended:

\begin{equation}
    q_{j}
    =
    \frac{r_{j}}{R_{L}},
\end{equation}

\begin{equation}
    w_{j}
    =
    \begin{cases}
        \left(1-q_{j}\right)^{2},
        & q_{j}<1,\\
        0,
        & q_{j}\geq 1.
    \end{cases}
\end{equation}

The light contribution to boundary sample $j$ is

\begin{equation}
    \Delta \mathbf{L}_{j}
    =
    w_{j}
    E_{\mathrm{deposit}}
    \mathbf{C}_{p},
\end{equation}

where $\mathbf{C}_{p}$ is the photon color.

The weights should preferably be normalized:

\begin{equation}
    \widehat{w}_{j}
    =
    \frac{w_{j}}
    {\sum_{k}w_{k}+\epsilon}.
\end{equation}

The normalized contribution becomes

\begin{equation}
    \Delta \mathbf{L}_{j}
    =
    \widehat{w}_{j}
    E_{\mathrm{deposit}}
    \mathbf{C}_{p}.
\end{equation}

Normalization prevents the total deposited light from increasing when more
boundary samples are found inside the influence radius.

\section{Finding Neighboring Boundary Samples}

The implementation needs a method for finding boundary samples near the photon
hit location.

The preferred method is to reuse the existing spatial cell structure.

The procedure is:

\begin{enumerate}
    \item Convert the photon hit position to a spatial cell index.
    \item Inspect the hit cell and immediately adjacent cells.
    \item Read the boundary samples stored in those cells.
    \item Calculate the distance from each candidate boundary sample to the hit.
    \item Apply the deposition kernel when the distance is less than $R_{L}$.
\end{enumerate}

Because the existing particle framework already uses grid-cell occupancy, the
light-deposition system should avoid introducing an independent neighborhood
search unless necessary.

\section{Atomic Accumulation}

Multiple photons may deposit light into the same boundary sample during one
dispatch. The accumulation operation must therefore be atomic.

Conceptually:

\begin{lstlisting}[language=C++]
atomicAdd(
    boundaryLightCurrent[boundaryIndex].red,
    depositedRed
);

atomicAdd(
    boundaryLightCurrent[boundaryIndex].green,
    depositedGreen
);

atomicAdd(
    boundaryLightCurrent[boundaryIndex].blue,
    depositedBlue
);
\end{lstlisting}

If shader atomic floating-point operations are supported, floating-point
accumulation may be used directly.

Otherwise, use fixed-point integers:

\begin{lstlisting}[language=C++]
uint encodedRed =
    uint(max(redContribution, 0.0) * LIGHT_SCALE);

atomicAdd(
    boundaryLightCurrent[boundaryIndex].red,
    encodedRed
);
\end{lstlisting}

The render or filtering pass converts the stored integer value back to a
floating-point value.

\section{Temporal Accumulation}

Even with spatial spreading, sparse photon samples may cause frame-to-frame
noise. The illumination should therefore be filtered over time.

Let

\begin{itemize}
    \item $\mathbf{L}_{j}^{n}$ be the previous filtered light,
    \item $\mathbf{D}_{j}^{n+1}$ be the current-frame deposited light,
    \item $\alpha$ be the current-frame contribution factor.
\end{itemize}

Use

\begin{equation}
    \mathbf{L}_{j}^{n+1}
    =
    \left(1-\alpha\right)
    \mathbf{L}_{j}^{n}
    +
    \alpha
    \mathbf{D}_{j}^{n+1}.
\end{equation}

For strong smoothing, use a small $\alpha$, such as

\begin{equation}
    0.02 \leq \alpha \leq 0.10.
\end{equation}

An alternative persistence form is

\begin{equation}
    \mathbf{L}_{j}^{n+1}
    =
    \lambda
    \mathbf{L}_{j}^{n}
    +
    \mathbf{D}_{j}^{n+1},
\end{equation}

where $\lambda$ is typically close to one.

For example,

\begin{equation}
    0.90 \leq \lambda \leq 0.99.
\end{equation}

The normalized exponential-moving-average form is generally easier to control
because the illumination does not grow indefinitely under constant lighting.

\section{Light-Buffer Normalization}

The current-frame deposited light depends on the number of emitted photon
samples. Increasing the photon count should not automatically make the physical
light source brighter unless explicitly intended.

The deposited illumination should therefore be normalized by the photon
sampling rate.

If the emitter represents a total light intensity $P_{e}$ and emits $N_{p}$
photons per sampling interval, assign each photon

\begin{equation}
    E_{p}
    =
    \frac{P_{e}}{N_{p}}.
\end{equation}

This allows the photon count to control sampling quality without changing total
emitter intensity.

\section{Boundary Rendering}

The boundary fragment shader should read the filtered illumination value for the
associated boundary sample.

A basic rendering equation is

\begin{equation}
    \mathbf{C}_{\mathrm{final}}
    =
    \mathbf{C}_{\mathrm{base}}
    \left(
        A_{\mathrm{ambient}}
        +
        G\mathbf{L}_{\mathrm{filtered}}
    \right),
\end{equation}

where

\begin{itemize}
    \item $\mathbf{C}_{\mathrm{base}}$ is the boundary base color,
    \item $A_{\mathrm{ambient}}$ is a small ambient term,
    \item $G$ is an exposure or gain factor,
    \item $\mathbf{L}_{\mathrm{filtered}}$ is accumulated photon illumination.
\end{itemize}

In GLSL:

\begin{lstlisting}[language=C++]
vec3 baseColor = boundaryMaterial.baseColor.rgb;
vec3 photonLight = boundaryLightFiltered[boundaryIndex].rgb;

vec3 finalColor =
    baseColor *
    (ambientLevel + photonLight * lightExposure);

outColor = vec4(finalColor, 1.0);
\end{lstlisting}

A tone-mapping operation may be applied when illumination values have a large
range:

\begin{equation}
    \mathbf{C}_{\mathrm{mapped}}
    =
    \frac{\mathbf{C}}
    {\mathbf{C}+\mathbf{1}}.
\end{equation}

In GLSL:

\begin{lstlisting}[language=C++]
finalColor = finalColor / (finalColor + vec3(1.0));
\end{lstlisting}

\section{Optional Directional Information}

If desired, each boundary-light sample may also accumulate the incoming photon
direction.

Add the following values:

\begin{lstlisting}[language=C++]
struct BoundaryLightDirectional
{
    glm::vec4 accumulatedLight;
    glm::vec4 accumulatedDirection;
};
\end{lstlisting}

For each photon hit:

\begin{equation}
    \Delta \mathbf{D}_{j}
    =
    \widehat{w}_{j}
    E_{\mathrm{deposit}}
    \left(
        -\mathbf{d}_{p}
    \right).
\end{equation}

The average incoming-light direction is

\begin{equation}
    \overline{\mathbf{D}}_{j}
    =
    \frac{\mathbf{D}_{j}}
    {\left\|\mathbf{D}_{j}\right\|+\epsilon}.
\end{equation}

This direction may be used for additional diffuse or specular surface shading.

However, directional accumulation is optional. The first implementation should
accumulate color and intensity only.

\section{Optional Photon Tracers}

Photon tracers may still be rendered to show the light transport paths.

They should not be responsible for illuminating the surface.

To reduce clutter and rendering cost, render only a subset of the simulated
photons.

For example:

\begin{lstlisting}[language=C++]
bool renderPhoton =
    (photonIndex % tracerSamplingInterval) == 0;
\end{lstlisting}

The tracer color may use the eye-alignment equation described previously.

A radial point-sprite falloff may be used:

\begin{lstlisting}[language=C++]
vec2 point = gl_PointCoord * 2.0 - 1.0;
float radiusSquared = dot(point, point);

if (radiusSquared > 1.0)
{
    discard;
}

float radialWeight = exp(-4.0 * radiusSquared);

vec3 toEye =
    normalize(cameraPosition - photonPosition);

float alignment =
    max(dot(photonDirection, toEye), 0.0);

float angularWeight =
    minimumTracerBrightness +
    (1.0 - minimumTracerBrightness) *
    pow(alignment, tracerExponent);

float alpha = radialWeight * angularWeight;
\end{lstlisting}

Additive blending may be used for the optional photon tracers.

\section{Recommended GPU Pass Order}

The recommended frame sequence is:

\begin{enumerate}
    \item Clear the current-frame boundary-light buffer.
    \item Update photon positions.
    \item Detect photon-boundary collisions.
    \item Deposit photon energy into nearby boundary samples.
    \item Apply photon reflection, termination, or rebirth.
    \item Insert a synchronization barrier.
    \item Temporally filter the current boundary illumination.
    \item Insert a synchronization barrier.
    \item Render boundaries using the filtered illumination buffer.
    \item Optionally render photon tracers.
\end{enumerate}

A conceptual pipeline is:

\begin{lstlisting}
Clear Current Light Buffer
            |
            v
Update Photon Positions
            |
            v
Detect Boundary Hits
            |
            v
Deposit Light into Boundary Samples
            |
            v
Temporal Filtering
            |
            v
Render Illuminated Boundaries
            |
            v
Render Optional Photon Tracers
\end{lstlisting}

\section{Synchronization Requirements}

The deposition pass writes the current boundary-light buffer.

The temporal filtering pass reads the current boundary-light buffer and writes
the filtered-light buffer.

The boundary rendering pass reads the filtered-light buffer.

Required dependencies are therefore:

\begin{enumerate}
    \item Deposition shader writes must be visible to the filtering shader.
    \item Filtering shader writes must be visible to the graphics shader.
\end{enumerate}

In Vulkan, appropriate compute-to-compute and compute-to-graphics barriers must
be inserted.

Conceptually:

\begin{lstlisting}[language=C++]
Photon deposition:
    shader write -> boundaryLightCurrentBuffer

Barrier:
    compute shader write
    to compute shader read

Temporal filtering:
    read boundaryLightCurrentBuffer
    write boundaryLightFilteredBuffer

Barrier:
    compute shader write
    to vertex/fragment shader read

Boundary rendering:
    read boundaryLightFilteredBuffer
\end{lstlisting}

The exact stage and access masks should match the shader stages that read and
write each buffer.

\section{Double Buffering}

If the temporal filter reads and writes the same illumination array, a race may
occur between workgroups.

The preferred solution is to use two filtered-light buffers:

\begin{lstlisting}[language=C++]
VkBuffer boundaryLightFilteredPrevious;
VkBuffer boundaryLightFilteredNext;
\end{lstlisting}

The temporal filter reads the previous buffer and writes the next buffer:

\begin{equation}
    \mathbf{L}_{\mathrm{next}}
    =
    \left(1-\alpha\right)
    \mathbf{L}_{\mathrm{previous}}
    +
    \alpha
    \mathbf{D}_{\mathrm{current}}.
\end{equation}

After the filtering dispatch, swap the two buffer handles.

This avoids in-place read-write hazards.

\section{Debugging Outputs}

The first implementation should include optional debug counters.

Recommended counters include:

\begin{itemize}
    \item number of active photons,
    \item number of boundary collisions,
    \item number of valid light deposits,
    \item number of photons terminated,
    \item number of photons reflected,
    \item total deposited fixed-point energy,
    \item maximum illumination value,
    \item number of failed boundary-neighbor searches.
\end{itemize}

A possible structure is:

\begin{lstlisting}[language=C++]
struct PhotonLightingDebug
{
    uint32_t activePhotonCount;
    uint32_t collisionCount;
    uint32_t depositCount;
    uint32_t terminatedCount;

    uint32_t reflectedCount;
    uint32_t failedNeighborSearchCount;
    uint32_t maximumEncodedLight;
    uint32_t padding;
};
\end{lstlisting}

These counters should be reset before each debug frame.

\section{Implementation Stages}

The implementation should be completed incrementally.

\subsection{Stage 1: Single-Sample Deposition}

Implement:

\begin{itemize}
    \item photon-boundary hit detection,
    \item incidence-angle calculation,
    \item atomic deposit into the directly hit boundary sample,
    \item boundary rendering using the light buffer.
\end{itemize}

Do not initially implement spatial smoothing or temporal filtering.

This stage verifies that light is deposited on the correct boundary and that the
direction sign is correct.

\subsection{Stage 2: Temporal Filtering}

Add:

\begin{itemize}
    \item previous and next filtered-light buffers,
    \item exponential moving average,
    \item filtered-light buffer swapping.
\end{itemize}

This stage should reduce frame-to-frame flicker.

\subsection{Stage 3: Spatial Spreading}

Add:

\begin{itemize}
    \item boundary-neighbor lookup,
    \item influence radius,
    \item compact smoothing kernel,
    \item normalized kernel weights.
\end{itemize}

This stage should remove isolated bright boundary samples.

\subsection{Stage 4: Material Response}

Add:

\begin{itemize}
    \item absorption,
    \item reflection,
    \item photon energy reduction,
    \item optional surface color filtering.
\end{itemize}

\subsection{Stage 5: Optional Photon Tracers}

Add a separate rendering path for a sampled subset of photons.

Tracer brightness may depend on eye alignment, but tracer brightness must not
affect deposited surface illumination.

\section{Initial Default Parameters}

The following values are recommended as initial tuning values:

\begin{lstlisting}[language=C++]
float photonSpeed = ...;              // existing fixed photon speed
float photonEnergy = 1.0f / photonCountPerEmitter;

float lightInfluenceRadius = 1.5f;    // relative to boundary spacing
float temporalAlpha = 0.05f;
float ambientLevel = 0.03f;
float lightExposure = 1.0f;

float tracerMinimumBrightness = 0.03f;
float tracerExponent = 3.0f;

uint32_t tracerSamplingInterval = 10;
\end{lstlisting}

These values are starting points and should be exposed through the simulation
configuration.

\section{Configuration Items}

The following configuration values should be added:

\begin{lstlisting}
photon_lighting_enabled = true;

photon_energy_normalization = true;
photon_light_exposure = 1.0;
photon_light_ambient = 0.03;

photon_light_temporal_filter = true;
photon_light_temporal_alpha = 0.05;

photon_light_spatial_filter = true;
photon_light_influence_radius = 1.5;

photon_tracers_enabled = true;
photon_tracer_sampling_interval = 10;
photon_tracer_minimum_brightness = 0.03;
photon_tracer_exponent = 3.0;
\end{lstlisting}

Optional material properties are:

\begin{lstlisting}
boundary_light_absorption = 1.0;
boundary_light_reflection = 0.0;
boundary_light_roughness = 0.0;
\end{lstlisting}

\section{Acceptance Criteria}

The implementation is considered successful when the following conditions are
met:

\begin{enumerate}
    \item Boundary illumination is visible even when individual photons are
    sparse.
    \item The illuminated surface is significantly smoother than the rendered
    photon distribution.
    \item Boundary brightness depends on photon incidence angle, not camera
    angle.
    \item Camera movement does not change the amount of deposited boundary
    light.
    \item Camera movement may change only the visibility of optional photon
    tracers.
    \item Increasing the number of photons improves sampling quality without
    proportionally increasing total light intensity.
    \item Temporal filtering reduces flicker.
    \item Spatial deposition reduces isolated bright boundary samples.
    \item Multiple photons may safely update the same boundary-light sample.
    \item Boundary rendering reads only the completed filtered-light buffer.
\end{enumerate}

\section{Summary of Required Code Changes}

Codex should implement the following changes:

\begin{enumerate}
    \item Add per-photon color and energy values if they do not already exist.
    \item Add current and filtered boundary-light storage buffers.
    \item Clear the current light buffer at the start of each lighting update.
    \item Extend photon-boundary collision handling to calculate the boundary
    normal and photon incidence factor.
    \item Deposit photon energy into the boundary-light buffer using atomic
    operations.
    \item Add spatial deposition over neighboring boundary samples.
    \item Add a temporally filtered double-buffered lighting pass.
    \item Add Vulkan barriers between deposition, filtering, and rendering.
    \item Modify the boundary rendering shaders to read accumulated photon
    illumination.
    \item Keep the eye-alignment calculation only in the optional photon-tracer
    rendering path.
    \item Normalize photon energy by the number of emitted photon samples.
    \item Add configuration controls and debug counters.
\end{enumerate}

The central design requirement is that photon particles transport and deposit
light, while the boundaries display the accumulated illumination. Direct
photon visibility is optional and must remain independent from surface lighting.

\end{document}
```
